\scanpage{37}
\margin{$?\quad\big|$}Une question plus subtile est s'ils peuvent
s'inverser, quand c'est des isomorphismes ....

Dualité entre $P^i(X)$ et $P^{\deleted{n+1-i}\inserted{j}}(X)$ \inserted{pour $i+j=n+1$} : Si $A,B$
\deleted{\unsure{sont une|sont deux}} \inserted{des} VA, \deleted{dont la dualité est} accouplées
par une dualité parfaite $(A\simeq B^{\unsure{*|\vee}}, B\simeq A^*)$,
alors il y a correspondance biunivoque
entre $\operatorname{Hom\,alg}(\deleted{\unsure{B}}A,P^i(X))$ et $\operatorname{Hom}^{\mathrm{alg}}(P^{\unsure{1|j}}(X),B)$ :

En effet, si on a un hom
$A\longrightarrow P^i(X)$, défini par un $Z\in A^i_{\mathrm{alb}}(\deleted{A}X\times\deleted{X}A)$,
alors $ {}^t Z$ définit un hom $P^j(X)\longrightarrow P^{\deleted{\unsure{1}}i+j-n}(A)=\underset{\substack{\shortparallel\\B}}{P^1(A)}$
et inversement si on a $P^j(Y)\longrightarrow B$. \deleted{// Notation}
\deleted{$\unsure{{}^t\varphi|\mathrm{tq}}$ \unsure{s'impose}. Plus généralement}
Cependant, \margin{$\big|\big|\big|$}il faut vérifier que $P^j(X)\longrightarrow B$ ne dépend que
de \uline{l'application} $A\longrightarrow P^i(X)$, pas du cycle $Z$.

Question de transposition se pose, plus généralement, pour
$P^i(X)\unsure{-|\longrightarrow}P^{\unsure{i'|a'|j'}}(X')$ en $P^{j'}(X')\unsure{-|\longrightarrow}P^j(X)$.
Cf plus bas ....
% archive sheet number: 36
